\section{Introduction}

Introduction to the problem and the approach.

\section{Related Work}

Overview of related work.

\section{Methodology}

A description of the data and methods.

\section{Results}

This section includes the results of the experiments, primarily in tabular or graphical form.

\section{Discussion}

Discussion of the results and their implications.

\section{Future Work}

What would you do next?

\section{Conclusions}

Summary of the work and its contributions.

\section*{Acknowledgements}

We thank the Data Science at Georgia Tech (DS@GT) CLEF competition group for their support.
This research was supported in part through research cyberinfrastructure resources and services provided by the Partnership for an Advanced Computing Environment (PACE) at the Georgia Institute of Technology, Atlanta, Georgia, USA \cite{PACE}. 

%% The declaration on generative AI comes in effect
%% in Janary 2025. See also
%% https://ceur-ws.org/GenAI/Policy.html
\section*{Declaration on Generative AI}
  {\em Either:}\newline
  The author(s) have not employed any Generative AI tools.
  \newline
  
 \noindent{\em Or (by using the activity taxonomy in ceur-ws.org/genai-tax.html):\newline}
 During the preparation of this work, the author(s) used X-GPT-4 and Gramby in order to: Grammar and spelling check. Further, the author(s) used X-AI-IMG for figures 3 and 4 in order to: Generate images. After using these tool(s)/service(s), the author(s) reviewed and edited the content as needed and take(s) full responsibility for the publication’s content. 